\documentclass[10pt,a4paper]{amsart}
\usepackage[margin=19mm]{geometry}
\usepackage{amsmath,amssymb,mathtools}
\usepackage[scr=rsfs,cal=euler]{mathalfa}
\usepackage{xfrac}
\usepackage[unicode,hidelinks,bookmarks=false]{hyperref}
\newcommand{\ie}{i.e.\ }
\newcommand{\cf}{cf.\ }
\newcommand{\resp}{respectively\ }
\newcommand{\Subsection}{\subsection}
\providecommand{\nobreakdash}{\nobreak\hbox{-}}
\setlength{\emergencystretch}{2em}
\multlinegap0pt
\usepackage{bm}
\usepackage{bbm}
\usepackage{dsfont}
\usepackage{xspace}
\usepackage{cancel}
\usepackage{mathtools}
\usepackage{xcolor}
\usepackage{amsthm}
\makeatletter
\@ifundefined{definition}{\newtheorem{definition}{Definition}}{}
\@ifundefined{lemma}{\newtheorem{lemma}[definition]{Lemma}}{}
\makeatother
\title[A nonlocal Aw-Rascle-Zhang system with linear pressure term]{A nonlocal Aw-Rascle-Zhang system with linear pressure term}
\author{Debora Amadori, Felisia Angela Chiarello, Gianmarco Cipollone}
\date{Source: arXiv:2509.20973v3 (CC BY 4.0)}
\begin{document}
\maketitle

\noindent\emph{Source and licence.} A nonlocal Aw-Rascle-Zhang system with linear pressure term. Version arXiv:2509.20973v3 (CC BY 4.0), distributed under the Creative Commons Attribution 4.0 International licence, as recorded in the arXiv licence metadata for this exact version and shown on the abstract page https://arxiv.org/abs/2509.20973v3. Full rights notice: licenses/arxiv-2509.20973-CC-BY-4.0.txt.\par\smallskip

\noindent\textit{Editorial scope.} Counted unit: the Lemma whose source label is \texttt{lemma: barycentric} in the article (source0.tex lines 447--455, source label \texttt{lemma: barycentric}), delivered together with the article's own complete original proof of it (source0.tex lines 457--511, one \texttt{proof} environment of 55 lines). No mathematics is written, repaired or re-derived: statement and proof are the article's own text line for line. Required context retained uncounted: def:i-star (lines 286--288). Every definition, display and label of those lines is retained unchanged. Omitted from this package and not redistributed: 1--285, 289--446, 456--456, 512--1652. Nothing inside a retained excerpt is omitted or altered: every retained body is delivered exactly as the article's lines are written, so no omission receipt of any kind accompanies this package. Citations: the retained excerpts carry no citation command at all, so no bibliography entry of the article is transcribed and no reference is redistributed. Reference adaptation: none was needed and none was made; every reference inside the delivered unit resolves inside it, so no unresolved reference or citation is produced. Printed slips kept literal: no printed word, symbol, hypothesis or number of the counted statement or of its proof was altered. Layout and class: the article's own class is \texttt{amsart} with packages including inputenc, fontenc, graphicx, babel, bold-extra, geometry, algorithm, amsmath, mathtools, microtype, url, cite, listings, amsaddr; the wrapper is the lane's pinned \texttt{amsart} editorial wrapper at 10pt on A4, which restates the theorem-like environments the retained text needs and adds the article's own macro definitions that the retained text needs (none), with extra wrapper packages (bm, bbm, dsfont, xspace, cancel, mathtools, xcolor). The delivered numbering is the unit's own. Assets: no retained span contains a figure environment, a graphics-inclusion command, a PDF-page inclusion, a table or a TikZ picture; no image file is copied into this package, the package declares no asset dependency and no image file is redistributed with it. Licence: version arXiv:2509.20973v3 of this article is distributed under the Creative Commons Attribution 4.0 International licence, as recorded in the arXiv licence metadata for this exact version and shown on the abstract page https://arxiv.org/abs/2509.20973v3. A full rights notice is delivered at licenses/arxiv-2509.20973-CC-BY-4.0.txt. Characters: every retained excerpt is pure ASCII, so the delivered engine, XeLaTeX with its default Latin Modern text font, reports no missing character.
\begin{equation}\label{def:i-star}
    i_*(t)=\min J_i(t),\qquad i^*(t)=\max J_i(t).
\end{equation}

\begin{lemma}[\textbf{The barycentric lemma}]\label{lemma: barycentric}
    Fix an $i\in\{1,\dots,N\}$ and a time $t>0$. For any $k\in J_i(t)$ the following holds
    \begin{equation}\label{barycentric equation}
        \frac{\sum_{j=i_*(t)}^km_j\psi_j(t-)}{\sum_{j=i_*(t)}^km_j}\geq
                \frac{\sum_{j\in J_i(t)}m_j\psi_j(t-)} {\sum_
          {{\color{black}j\in J_i(t)}} m_j}=\psi_i(t+)\geq\frac{\sum_{j=k}^{i^*(t)}m_j\psi_j(t-)}{\sum_{j=k}^{i^*(t)}m_j}
    \end{equation}
    {\color{black} where $i_*(t)$, $i^*(t)$ are defined in \eqref{def:i-star}.}
\end{lemma}

\begin{proof}
To prove the lemma it is sufficient to establish the following monotonicity property
    \begin{equation}\label{monotonicity of psi}
        \psi_{i_*(t)}(t-)\geq\dots\geq\psi_{i^*(t)}(t-).
    \end{equation}
     Let us show how we obtain \eqref{barycentric equation} from \eqref{monotonicity of psi}.\newline
    Without loss of generality we suppose to have some quantities $b_i$ and $\alpha_i$ for $i=1,\dots,M$ such that
    \begin{equation}\label{eq:monotonicity_b_i}
       b_1\ge b_2\ge\dots\ge b_M, \quad\sum_{i=1}^M\alpha_i=C,\quad C\in(0,1]\,.
    \end{equation}   
    We only prove the first inequality in \eqref{barycentric equation}, the proof is similar for the second one. 
    We consider a $\bar k\in\{1,\dots,M\}$ and the aim is to show that
    \begin{equation*}
        \frac{\sum_{j=1}^{\bar k}\alpha_jb_j}{\sum_{j=1}^{\bar k}\alpha_j}\geq\frac{\sum_{j=1}^M\alpha_jb_j}
        {\sum_{j=1}^M\alpha_j}=\frac{\sum_{j=1}^M\alpha_jb_j}{C},
    \end{equation*}
    that is equivalent to
    \begin{equation*}
        C\sum_{j=1}^{\bar k}\alpha_jb_j\geq\left(\sum_{j=1}^{\bar k}\alpha_j\right)\left(\sum_{j=1}^M\alpha_jb_j\right)
        =\left(\sum_{j=1}^{\bar k}\alpha_j\right)\left(\sum_{j=1}^{\bar k}\alpha_jb_j+\sum_{j=\bar k+1}^M\alpha_jb_j\right),
    \end{equation*}
    and this is true if and only if the following inequality holds true
    \begin{equation*}
        \left(C-\sum_{j=1}^{\bar k}\alpha_j\right)\left(\sum_{j=1}^{\bar k}\alpha_jb_j \right)\geq 
        \left(\sum_{j=1}^{\bar k}\alpha_j\right)\left(\sum_{j=\bar k+1}^M\alpha_jb_j\right).
    \end{equation*}
    We know that $\sum_{j=1}^M\alpha_j=C,$ then putting it in the last inequality we obtain the following
    \begin{equation*}
        \left(\sum_{j=1}^{\bar k}\alpha_jb_j\right)\left(\sum_{j=\bar k+1}^M\alpha_j\right)
        \geq\left(\sum_{j=1}^{\bar k}\alpha_j\right)\left(\sum_{j=\bar k+1}^M\alpha_jb_j\right).
    \end{equation*}
    Writing explicitly the sums, we have
    \begin{equation*}
        \left(\alpha_1b_1+\dots+\alpha_{\bar k}b_{\bar k}\right)\left(\sum_{j=\bar k+1}^M\alpha_j\right)
        \geq\left(\alpha_1+\dots+\alpha_{\bar k}\right)\left(\sum_{j=\bar k+1}^M\alpha_jb_j\right).
    \end{equation*}
    From \eqref{eq:monotonicity_b_i}  we obtain that the inequality holds true. 
    Substituting the $b_i$'s and the $\alpha_i$'s respectively with the $\psi_i(t-)$'s and the $m_i$'s, 
    we get \eqref{barycentric equation} if \eqref{monotonicity of psi} holds true.
    \newline
    It remains to prove \eqref{monotonicity of psi}, i.e. at a collision time $t$ one has
 \begin{equation*}
        \psi_{k}(t-)\geq\psi_{k+1}(t-)\qquad \forall\, k = i_*(t), \ldots, i^*(t)-1\,.
    \end{equation*}
    We know that for any $j\in J_i(t)$ we have
    $$
    \psi_k(t-)=v_k(t-)+\sum_{j=1}^N m_j\omega(x_k(t)-x_j(t)).
    $$
    Then, it is possible to conclude that \eqref{monotonicity of psi} follows 
    from the obvious monotonicity property of the velocities, i.e.
    \begin{equation*}
        v_{i_*(t)}(t-)\geq\dots\geq v_{i^*(t)}(t-),
    \end{equation*}
   that is true due to the fact that $t$ is a collision time. This completes the proof.
\end{proof}

\end{document}
